\documentclass[../main.tex]{subfiles}

\graphicspath{{../image/}{../figure/}}

\begin{document}

\chapter{连续介质物理导引}\label{chapter:ch02-连续介质物理导引}

在连续介质物理中，物体被建模为连续介质，其运动与平衡由平衡律和本构关系所支配。

平衡律的清单标识出理论，例如力学、热力学连续介质力学、电动力学等等。下面将给出典型平衡律的参考（Lagrange）与空间（Euler）表述，并记录划定连续介质热力学力学范围的质量、动量、能量平衡律以及 Clausius--Duhem 不等式。

本构关系的类型编码了材料响应的本质。下面将引入热弹性和热粘弹性的本构方程，并讨论热力学第二定律、材料标架无差异原理和材料对称性所施加的限制。

\subfile{Chapter-02/section-01.tex}
\subfile{Chapter-02/section-02.tex}
\subfile{Chapter-02/section-03.tex}
\subfile{Chapter-02/section-04.tex}
\subfile{Chapter-02/section-05.tex}
\subfile{Chapter-02/section-06.tex}
\subfile{Chapter-02/section-07.tex}
\subfile{Chapter-02/section-08.tex}
\subfile{Chapter-02/section-09.tex}

\end{document}
